\documentclass[11pt]{article}
\usepackage[utf8]{inputenc}
\usepackage{amsmath}
\usepackage{geometry}
\usepackage{listings}
\usepackage{xcolor}

\geometry{letterpaper, margin=1in}

\definecolor{codegreen}{rgb}{0,0.6,0}
\definecolor{backcolour}{rgb}{0.95,0.95,0.92}
\definecolor{codepurple}{rgb}{0.58,0,0.82}
\definecolor{codegray}{rgb}{0.5,0.5,0.5}

\lstdefinestyle{mystyle}{
    backgroundcolor=\color{backcolour},   
    commentstyle=\color{codegreen},
    keywordstyle=\color{blue},
    numberstyle=\tiny\color{codegray},
    stringstyle=\color{codepurple},
    basicstyle=\ttfamily\footnotesize,
    breaklines=true,                 
    numbers=left,                    
    numbersep=5pt,                  
    tabsize=2,
    showstringspaces=false
}

\lstset{style=mystyle}

\title{\textbf{Editorial: Problem G - The Discipline Committee's Struggle}}
\author{Setter: Hudson}
\date{}

\begin{document}

\maketitle

\section*{Problem Overview}
We need to support Range Add, Range Min (ChMin), Range Sum, and \textbf{Range Historic Sum} on an array.

\section*{Difficulty Analysis}
This problem combines two very advanced Segment Tree techniques:
\begin{enumerate}
    \item \textbf{Segment Tree Beats:} Required to handle Range ChMin updates efficiently.
    \item \textbf{Historic Information Maintenance:} Required to track the sum of values over time. 
\end{enumerate}
Because the "ChMin" operation selectively updates only values strictly greater than $x$, the history accumulation becomes conditional. This requires maintaining separate lazy tags for the maximum elements and the non-maximum elements.

\section*{Solution Approach}

\subsection*{1. Segment Tree Beats Logic}
We maintain $max1$ (maximum value), $max2$ (strict second maximum), and $cnt$ (count of max values) in each node.
For an update $A_i = \min(A_i, x)$:
\begin{itemize}
    \item If $x \ge max1$, ignore.
    \item If $max2 < x < max1$, update $max1$ to $x$. This is effectively adding $(x - max1)$ to all maximum elements.
    \item If $x \le max2$, recurse.
\end{itemize}

\subsection*{2. Adding History}
We treat the operations as separate affine transformations for:
\begin{itemize}
    \item Class A: The elements equal to $max1$.
    \item Class B: The elements strictly less than $max1$.
\end{itemize}
We maintain four lazy tags:
\begin{itemize}
    \item $lazy\_A$: Value added to maximums.
    \item $lazy\_B$: Value added to non-maximums.
    \item $hist\_A$: Historic value accumulated by maximums.
    \item $hist\_B$: Historic value accumulated by non-maximums.
\end{itemize}

When pushing tags down to a child:
\begin{itemize}
    \item If child's $max1$ matches parent's old $max1$, it inherits Class A tags.
    \item Otherwise, it inherits Class B tags.
\end{itemize}

\section*{Implementation}
The complexity is $O(Q \log^2 N)$ or $O(Q \log N)$ depending on the potential function analysis, but practically it is fast enough for $N=2 \cdot 10^5$.

\end{document}
